\chapter{TINJAUAN PUSTAKA}
\section{Penelitian Terdahulu}
Sintesis penelitian relevan dan jelaskan posisi penelitian yang dilakukan, bukan hanya rangkuman artikel.
\section{Landasan Teori}
Definisikan istilah, konsep, dan hasil matematika yang benar-benar digunakan.
\begin{definition}
Misalkan $x\in\mathbb{R}$. Definisi ini merupakan contoh lingkungan pernyataan matematis yang perlu diganti sesuai topik penelitian.
\end{definition}
Contoh persamaan bernomor:
\begin{equation}\label{eq:contoh}
 f(x)=x^2.
\end{equation}
Persamaan~\eqref{eq:contoh} hanya merupakan contoh penulisan.
